% ============================================================
% MÓDULO: MERCADOS FINANCIEROS Y BURSÁTILES
% ============================================================
\clearpage
\section*{MÓDULO: MERCADOS FINANCIEROS Y BURSÁTILES}
\subsection*{Bolsas, operaciones, brokers, índices y bancos centrales}

\begin{recuadro}
\textbf{Objetivo del módulo:} entender qué son los mercados de capitales, cómo funcionan las
bolsas argentinas y del mundo, cómo se opera desde un broker online, qué tipos de operaciones
existen (acciones, bonos, obligaciones negociables, cauciones, dólar MEP, futuros, opciones
\textit{call} y \textit{put}, venta en corto) y qué significan los indicadores que aparecen
todos los días en las noticias (Merval, riesgo país, Dow Jones, Nasdaq, tasa de la Fed).

\textbf{Importante:} este módulo es de \textbf{educación financiera}, no de recomendación de
inversiones. Las empresas, plataformas y activos que se nombran son ejemplos. Toda inversión
tiene riesgo y los resultados pasados no garantizan resultados futuros.
\end{recuadro}

% ============================================================
\section*{1. ¿QUÉ ES EL MERCADO DE CAPITALES?}

Es el ámbito donde se encuentran quienes \textbf{necesitan financiamiento} (empresas, gobiernos,
provincias) y quienes \textbf{tienen ahorros} y buscan hacerlos rendir (personas, fondos,
empresas, aseguradoras).

\begin{itemize}
  \item Una \textbf{empresa} que quiere crecer puede \textbf{emitir acciones} (vende una parte
        de la empresa a nuevos socios) o \textbf{emitir deuda} (obligaciones negociables).
  \item Un \textbf{gobierno} que necesita fondos \textbf{emite bonos o letras} (toma deuda).
  \item El \textbf{ahorrista} compra esos títulos esperando una ganancia: dividendos, intereses
        o la suba del precio.
\end{itemize}

\begin{center}
\begin{tikzpicture}[every node/.style={font=\small, align=center}, caja/.style={draw=ppAzul,
  very thick, rounded corners, minimum width=3.6cm, minimum height=1.2cm}]
  \node[caja] (A) at (0,0) {\textbf{AHORRISTAS}\\ personas, fondos};
  \node[caja] (M) at (5,0) {\textbf{MERCADO}\\ bolsa + brokers};
  \node[caja] (E) at (10,0) {\textbf{EMISORES}\\ empresas, Estado};
  \draw[->, thick] (A.north east) to[bend left=15] node[above]{\footnotesize dinero} (M.north west);
  \draw[->, thick] (M.north east) to[bend left=15] node[above]{\footnotesize dinero} (E.north west);
  \draw[->, thick] (E.south west) to[bend left=15] node[below]{\footnotesize títulos} (M.south east);
  \draw[->, thick] (M.south west) to[bend left=15] node[below]{\footnotesize títulos} (A.south east);
\end{tikzpicture}
\end{center}

\paragraph*{Mercado primario y secundario.}
En el \textbf{primario} el título se vende por primera vez y el dinero va a quien lo emite
(licitación de bonos del Tesoro, colocación de una obligación negociable, salida a bolsa de una
empresa u \textbf{IPO}). En el \textbf{secundario} los inversores se compran y venden los
títulos entre sí: eso es lo que ocurre todos los días en la bolsa, y es lo que le da
\textbf{liquidez} al primario (nadie compraría un bono a 10 años si no pudiera venderlo antes).

\paragraph*{Renta fija y renta variable.}
\begin{itemize}
  \item \textbf{Renta variable (acciones):} sos socio de la empresa. Ganás si sube el precio o
        si reparte dividendos; podés perder si baja. Mayor riesgo y mayor potencial.
  \item \textbf{Renta fija (bonos, obligaciones negociables, letras):} le prestás dinero a un
        gobierno o empresa, que se compromete a devolverlo con intereses en fechas
        determinadas. En general menor riesgo, aunque depende de quién emite: un bono de un
        país que no paga puede perder la mayor parte de su valor.
\end{itemize}

\paragraph*{Riesgo y rendimiento.}
Regla básica: \textbf{a mayor rendimiento esperado, mayor riesgo}. Nadie ofrece ganancias altas
y seguras; si alguien lo hace, probablemente sea una estafa. Para reducir el riesgo se
\textbf{diversifica}: no poner todos los huevos en la misma canasta.

% ============================================================
\section*{2. LOS MERCADOS ARGENTINOS}

\subsubsection*{2.1. Los actores}
\begin{center}
\small
\begin{tabular}{|p{0.27\linewidth}|p{0.63\linewidth}|}
\hline
\textbf{Actor} & \textbf{Función} \\ \hline
\textbf{CNV} (Comisión Nacional de Valores) & Organismo del Estado que regula y controla el
mercado de capitales: autoriza emisiones, agentes y fondos, y protege al inversor. \\ \hline
\textbf{BYMA} (Bolsas y Mercados Argentinos) & Principal mercado de acciones, bonos, cauciones,
opciones y CEDEARs del país. Calcula el índice S\&P Merval. \\ \hline
\textbf{A3 Mercados} & Nace en 2025 de la fusión del MAE (Mercado Abierto Electrónico, renta
fija y divisas) y Matba Rofex (futuros agropecuarios y financieros, con sede en Rosario). Es
el principal mercado de futuros y opciones del país (dólar futuro, soja, maíz, trigo). \\ \hline
\textbf{MAV} (Mercado Argentino de Valores) & Mercado con sede en Rosario, orientado a pymes
(cheques de pago diferido, pagarés, obligaciones negociables pyme). \\ \hline
\textbf{Bolsa de Comercio de Rosario (BCR)} & Institución centenaria vinculada al comercio de
granos: su \textbf{Cámara Arbitral} fija los precios de referencia (``pizarra'') de la soja,
el maíz y el trigo, que usa todo el agro argentino. \\ \hline
\textbf{Caja de Valores} & Custodia los títulos y registra quién es el dueño de cada uno. \\
\hline
\textbf{ALyC y otros agentes} & Los \textbf{brokers} (antes ``sociedades de bolsa''): abren la
cuenta del inversor y operan en su nombre. Deben estar inscriptos en la CNV. \\ \hline
\end{tabular}
\end{center}

\subsubsection*{2.2. Cómo funciona una rueda de negociación}
\begin{itemize}
  \item \textbf{La rueda:} es la jornada de negociación. En BYMA se opera de lunes a viernes
        hábiles, aproximadamente entre las 10:30 y las 17:00 h (desde noviembre de 2026, las
        acciones y CEDEARs a 24 h se pueden operar hasta las 18:00 h).
  \item \textbf{Ticker o especie:} código abreviado de cada título. Ejemplos: \texttt{GGAL}
        (Grupo Galicia), \texttt{YPFD} (YPF), \texttt{AL30} (bono Bonar 2030).
  \item \textbf{Moneda:} muchos títulos cotizan en pesos y también en dólares. El sufijo
        \textbf{D} indica dólar MEP (\texttt{AL30D}) y el sufijo \textbf{C} indica dólar cable
        o CCL (\texttt{AL30C}).
  \item \textbf{Plazos de liquidación:} cuándo se entregan efectivamente los títulos y el dinero.
        \textbf{Contado inmediato (CI o T+0):} el mismo día. \textbf{24 horas (T+1):} el día hábil
        siguiente; es el plazo estándar.
  \item \textbf{Puntas:} la \textbf{punta compradora} es el mayor precio que alguien ofrece
        pagar; la \textbf{punta vendedora}, el menor precio al que alguien acepta vender. La
        diferencia se llama \textbf{spread}: cuanto más chico, más líquido es el título.
  \item \textbf{Volumen:} cantidad de dinero o de títulos negociados. Un título con poco volumen
        puede ser difícil de vender rápido sin bajar el precio.
  \item \textbf{Paneles:} el \textbf{panel líder} reúne las acciones más negociadas; el
        \textbf{panel general}, el resto. También hay paneles de bonos, CEDEARs y pymes.
\end{itemize}

\subsubsection*{2.3. El índice S\&P Merval}
Un \textbf{índice bursátil} es un número que resume cómo se mueve un grupo de acciones. Es como
un ``termómetro'' del mercado: si el índice sube, la mayoría de las acciones que lo componen
subieron.

El \textbf{S\&P Merval} reúne las acciones de empresas argentinas más negociadas en BYMA
(bancos, energéticas, de servicios públicos, siderúrgicas, de telecomunicaciones). Como se
expresa en pesos, en una economía con inflación sube mucho en números; por eso los analistas
suelen mirarlo \textbf{en dólares} para saber si realmente ganó valor.

\subsubsection*{2.4. El riesgo país}
El \textbf{riesgo país} mide cuánto \textbf{más interés} le exigen los inversores a un país
para prestarle, en comparación con el gobierno de Estados Unidos (considerado el deudor más
seguro). Lo calcula el banco JP Morgan con su índice \textbf{EMBI} y se expresa en
\textbf{puntos básicos} (100 puntos básicos = 1\%).

\begin{recuadro}
\textbf{Ejemplo:} si los bonos del Tesoro de EE.\,UU. rinden 4\% anual y el riesgo país es de
850 puntos básicos, los bonos argentinos en dólares rinden aproximadamente
4\% + 8,5\% = \textbf{12,5\% anual}. Cuanto más alto el riesgo país, más caro le resulta al
Estado (y a las empresas y provincias) endeudarse, y más difícil es conseguir crédito
internacional.

\textbf{Dato:} a mediados de noviembre de 2026 el riesgo país argentino rondaba los 850 puntos
básicos, después de haber tocado un mínimo de 402 puntos en julio. Consultar el valor del día.
\end{recuadro}

\subsubsection*{2.5. Los bonos argentinos}
\begin{itemize}
  \item \textbf{Bonares (AL) y Globales (GD):} bonos del Tesoro en dólares. Los Bonares se rigen
        por ley argentina; los Globales, por ley de Nueva York (en caso de conflicto, se
        litiga en tribunales de EE.\,UU.). Ejemplos: AL30, GD30, GD35, GD46 (el número indica
        el año de vencimiento).
  \item \textbf{Letras y bonos capitalizables en pesos (Lecap, Boncap):} deuda en pesos a tasa
        fija, de corto plazo.
  \item \textbf{Bonos ajustables por inflación (CER):} el capital se actualiza con la inflación.
  \item \textbf{Bonos \textit{dollar linked}:} en pesos, pero ajustados por el tipo de cambio
        oficial; sirven para cubrirse de una devaluación.
  \item \textbf{Bonos del BCRA (BOPREAL)} y \textbf{bonos provinciales} (Provincia de Buenos
        Aires, Córdoba, Neuquén, etc.).
\end{itemize}

\paragraph*{Cómo se lee un bono.}
\begin{itemize}
  \item \textbf{Valor nominal (VN):} el monto sobre el que se calculan intereses y amortización.
        Los precios suelen expresarse cada 100 de valor nominal.
  \item \textbf{Cupón:} el interés que paga, en fechas fijas (por ejemplo, cada seis meses).
  \item \textbf{Amortización:} la devolución del capital, en una o varias cuotas.
  \item \textbf{Vencimiento:} fecha del último pago.
  \item \textbf{Paridad:} precio dividido valor técnico; una paridad de 70\% significa que el
        mercado paga 70 por algo que, si se cumple, vale 100.
  \item \textbf{TIR:} lo que rinde por año si se compra hoy y se mantiene hasta el vencimiento.
        \textbf{Cuando el precio del bono baja, su TIR sube}, y viceversa.
\end{itemize}

\subsubsection*{2.6. Las obligaciones negociables (ON)}
Son \textbf{bonos emitidos por empresas}. Con ellas las compañías financian inversiones
(perforar pozos, comprar equipos, construir plantas) sin pedir préstamos al banco. Hay ON de
grandes empresas (YPF, Pampa Energía, Telecom, Tecpetrol) y ON pyme.
\begin{itemize}
  \item Se compran en la \textbf{licitación primaria} (el broker muestra las ofertas abiertas y
        se ofertan montos y tasas) o después en el mercado secundario.
  \item Muchas están en \textbf{dólares} y pagan intereses periódicos; suelen tener montos
        mínimos de compra.
  \item El riesgo depende de la empresa: hay que mirar su \textbf{calificación} (rating) y su
        situación financiera.
\end{itemize}

% ============================================================
\section*{3. BROKERS: CÓMO SE INVIERTE ONLINE EN ARGENTINA}

\subsubsection*{3.1. Los agentes registrados en la CNV}
\begin{center}
\small
\begin{tabular}{|p{0.30\linewidth}|p{0.60\linewidth}|}
\hline
\textbf{Tipo de agente} & \textbf{Qué hace} \\ \hline
\textbf{ALyC} (Agente de Liquidación y Compensación) & Opera en los mercados, recibe el dinero
y los títulos de los clientes y liquida las operaciones. Es el ``broker'' propiamente dicho. \\
\hline
\textbf{AN} (Agente de Negociación) & Opera en el mercado pero liquida a través de un ALyC. \\
\hline
\textbf{AAGI} (Agente Asesor Global de Inversión) & Asesora y administra carteras de clientes. \\
\hline
\textbf{AP} (Agente Productor) & Consigue clientes y los deriva a un ALyC (comercializador). \\
\hline
\textbf{Sociedades Gerentes de FCI} & Administran fondos comunes de inversión. \\ \hline
\end{tabular}
\end{center}

\subsubsection*{3.2. Dónde se puede invertir}
\begin{itemize}
  \item \textbf{Bancos:} desde el homebanking, con la \textbf{cuenta títulos} asociada a la caja
        de ahorro. Cómodo, pero suelen tener menos instrumentos y comisiones más altas.
  \item \textbf{Brokers online (ALyC):} plataformas web y aplicaciones como Bull Market Brokers,
        InvertirOnline (IOL), Balanz, Cocos Capital, Portfolio Personal Inversiones (PPI) o
        Allaria, entre muchas otras. Ofrecen todo el mercado (acciones, bonos, ON, CEDEARs,
        cauciones, opciones, FCI), información, gráficos y capacitación.
  \item \textbf{Billeteras virtuales:} muchas permiten invertir en FCI money market, comprar
        dólar MEP o algunos instrumentos, a través de un ALyC asociado.
\end{itemize}
\textbf{Antes de elegir}, verificar en el sitio de la CNV que el agente esté registrado y
comparar \textbf{comisiones}, \textbf{costos de mantenimiento}, instrumentos disponibles,
atención al cliente y facilidad de la plataforma.

\subsubsection*{3.3. Paso a paso para empezar}
\begin{enumerate}
  \item \textbf{Abrir la cuenta comitente} online con DNI, CUIL, datos personales y
        laborales, y una selfie de validación. Es gratuita en la mayoría de los brokers.
  \item \textbf{Completar el test de perfil de inversor:} \textit{conservador} (prioriza no
        perder), \textit{moderado} o \textit{agresivo} (acepta más riesgo buscando más
        ganancia). El broker solo puede ofrecer productos acordes al perfil.
  \item \textbf{Vincular una cuenta bancaria propia} en pesos y en dólares. El dinero
        \textbf{solo puede entrar y salir desde cuentas a tu nombre}.
  \item \textbf{Transferir fondos} desde el banco o billetera a la cuenta del broker.
  \item \textbf{Operar:} buscar el ticker, elegir compra o venta, moneda, plazo, cantidad y
        tipo de orden; revisar el costo total y confirmar.
  \item \textbf{Controlar la tenencia} (valuación de la cartera) y los movimientos; descargar
        los comprobantes (boletos).
  \item \textbf{Retirar:} vender y transferir el dinero de vuelta al banco.
\end{enumerate}

\paragraph*{Los costos de operar.}
Cada operación tiene una \textbf{comisión del broker} (un porcentaje del monto), los
\textbf{derechos de mercado} que cobra la bolsa y el \textbf{IVA} sobre esos conceptos. En
operaciones chicas los costos pesan mucho: comprar y vender varias veces por semana con poco
dinero puede hacer perder plata aunque los precios suban.

% ============================================================
\section*{4. TIPOS DE ÓRDENES Y DE OPERACIONES}

\subsubsection*{4.1. Tipos de órdenes}
\begin{center}
\small
\begin{tabular}{|p{0.24\linewidth}|p{0.66\linewidth}|}
\hline
\textbf{Orden} & \textbf{Cómo funciona} \\ \hline
\textbf{A mercado} & Se ejecuta enseguida al mejor precio disponible. Rápida, pero el precio
final puede ser algo distinto al esperado. \\ \hline
\textbf{Límite} & Se fija el precio máximo de compra o mínimo de venta. Solo se ejecuta si el
mercado llega a ese precio. \\ \hline
\textbf{\textit{Stop loss}} & Orden de venta automática si el precio cae a un nivel fijado, para
limitar pérdidas. \\ \hline
\textbf{\textit{Take profit}} & Orden de venta automática si el precio sube a un nivel fijado,
para asegurar ganancias. \\ \hline
\end{tabular}
\end{center}

\subsubsection*{4.2. Compra y venta de acciones}
Se compra una acción esperando que suba o que pague \textbf{dividendos} (parte de las ganancias
que la empresa reparte entre sus accionistas, en efectivo o en nuevas acciones). Otros eventos
de una acción:
\begin{itemize}
  \item \textbf{Suscripción o ampliación de capital:} la empresa emite nuevas acciones y los
        accionistas tienen preferencia para comprarlas.
  \item \textbf{Split (desdoblamiento):} la acción se divide (por ejemplo, 1 acción de \$10.000
        pasa a ser 10 de \$1.000). No cambia el valor total, solo la hace más ``barata'' por
        unidad.
  \item \textbf{Asamblea:} los accionistas votan las decisiones importantes de la empresa.
\end{itemize}

\subsubsection*{4.3. Compra y venta de bonos}
Se puede comprar un bono para \textbf{cobrar sus cupones y amortizaciones hasta el vencimiento}
(se gana la TIR del momento de compra) o para \textbf{venderlo antes} si su precio sube (por
ejemplo, si baja el riesgo país). También se puede participar de las \textbf{licitaciones del
Tesoro}, donde el Estado coloca letras y bonos nuevos.

\subsubsection*{4.4. Caución bursátil y pases}
La \textbf{caución} es un \textbf{préstamo de dinero de muy corto plazo} (de 1 a 120 días)
entre inversores, garantizado con títulos y con la bolsa como garante.
\begin{itemize}
  \item \textbf{Caución colocadora:} prestás tus pesos (o dólares) a una tasa fija por unos días.
        Es una alternativa al plazo fijo, de bajo riesgo.
  \item \textbf{Caución tomadora:} pedís dinero prestado dejando tus títulos en garantía.
\end{itemize}
\textit{Ejemplo:} colocar \$1.000.000 a 7 días a una TNA del 30\% genera
$1.000.000 \times 0{,}30 \times 7/365 \approx$ \textbf{\$5.753} de interés (antes de costos).

El \textbf{pase} es una operación similar: se vende un título hoy y se pacta recomprarlo en una
fecha futura a un precio mayor (es usada por bancos y por el BCRA para manejar la liquidez).

\subsubsection*{4.5. Cheques de pago diferido y pagarés}
Una pyme que recibió un cheque a 90 días puede \textbf{venderlo en el mercado} (MAV, BYMA) y
cobrar hoy, con un descuento. Para el inversor que lo compra es una colocación de corto plazo
con una tasa mayor, avalada muchas veces por una SGR. Es una de las principales formas de
financiamiento de las pymes en el mercado de capitales.

\subsubsection*{4.6. Dólar MEP y dólar cable (CCL)}
El \textbf{dólar MEP} (o dólar bolsa) se obtiene \textbf{comprando un bono en pesos y
vendiéndolo en dólares}. El cociente entre ambos precios es el tipo de cambio.

\begin{recuadro}
\textbf{Ejemplo:} se compran 100 nominales de \texttt{AL30} en pesos a \$95.000 y se venden los
mismos 100 nominales como \texttt{AL30D} a USD 62.

Dólar MEP $= \dfrac{95.000}{62} \approx$ \textbf{\$1.532 por dólar}.

Los dólares quedan en la cuenta del broker y se transfieren a la caja de ahorro en dólares
propia. Hoy las personas humanas pueden hacerlo sin plazo de espera (\textit{parking}); las
personas jurídicas deben esperar un día hábil.
\end{recuadro}

El \textbf{dólar contado con liquidación (CCL o cable)} es igual, pero los dólares se reciben
en una cuenta \textbf{en el exterior}: se vende el bono como \texttt{AL30C}. Se usa, por
ejemplo, para comprar acciones en Wall Street o para que empresas giren dinero afuera.

\subsubsection*{4.7. Derivados: futuros}
Un \textbf{contrato de futuro} es un acuerdo para comprar o vender algo (dólar, soja, un índice)
\textbf{en una fecha futura a un precio fijado hoy}. Se negocian en A3 Mercados.
\begin{itemize}
  \item \textbf{Comprador del futuro (posición larga):} gana si el precio sube.
  \item \textbf{Vendedor del futuro (posición corta):} gana si el precio baja.
  \item No se paga el total al comienzo: se deposita una \textbf{garantía} (margen) y cada día
        se acreditan o debitan las diferencias (\textbf{ajuste diario}). Al vencimiento se
        liquida por diferencia en pesos.
\end{itemize}

\begin{recuadro}
\textbf{Ejemplo: dólar futuro.} Un taller que debe pagar en un mes una importación de USD 1.000
teme una devaluación. Compra un contrato de dólar futuro (USD 1.000) a fin de mes a
\textbf{\$1.600}.
\begin{itemize}[leftmargin=*,nosep]
  \item Si al vencimiento el dólar oficial está a \$1.650, el contrato le paga
        $(1.650 - 1.600) \times 1.000 =$ \textbf{\$50.000}, que compensan la suba del dólar.
  \item Si el dólar baja a \$1.550, pierde \$50.000 en el contrato, pero compra los dólares más
        baratos.
\end{itemize}
En los dos casos el taller ``fijó'' su dólar en \$1.600: eso es una \textbf{cobertura}. Un
productor de soja hace lo mismo para asegurar el precio de su cosecha antes de levantarla.
\end{recuadro}

\subsubsection*{4.8. Derivados: opciones \textit{call} y \textit{put}}
Una \textbf{opción} da el \textbf{derecho, pero no la obligación}, de comprar o vender una
acción a un precio fijado (\textbf{precio de ejercicio} o \textit{strike}) hasta una fecha de
vencimiento. Por ese derecho se paga una \textbf{prima}. En BYMA se operan sobre todo opciones
de Grupo Galicia (GGAL) y de algunas otras acciones líderes, en lotes de 100 acciones.

\begin{center}
\small
\begin{tabular}{|p{0.15\linewidth}|p{0.36\linewidth}|p{0.36\linewidth}|}
\hline
 & \textbf{CALL (opción de compra)} & \textbf{PUT (opción de venta)} \\ \hline
\textbf{Comprador} & Paga la prima. Tiene derecho a \textbf{comprar} al strike. Gana si la
acción \textbf{sube} por encima del strike + prima. & Paga la prima. Tiene derecho a
\textbf{vender} al strike. Gana si la acción \textbf{baja} por debajo del strike $-$ prima. \\
\hline
\textbf{Lanzador (vendedor)} & Cobra la prima. Está obligado a vender si el comprador ejerce. &
Cobra la prima. Está obligado a comprar si el comprador ejerce. \\ \hline
\textbf{Pérdida máxima del comprador} & La prima pagada. & La prima pagada. \\ \hline
\end{tabular}
\end{center}

\begin{recuadro}
\textbf{Ejemplo de CALL.} Una acción vale \$1.000. Compro un call con strike \$1.100 que vence en
dos meses, pagando una prima de \$50 por acción. Un lote son 100 acciones: invierto
\textbf{\$5.000}.
\begin{itemize}[leftmargin=*,nosep]
  \item Si al vencimiento la acción vale \$1.300: ejerzo, compro a \$1.100 lo que vale \$1.300.
        Ganancia $= (1.300 - 1.100 - 50) \times 100 =$ \textbf{\$15.000} (300\% de lo
        invertido).
  \item Si la acción queda en \$1.100 o menos: no ejerzo y pierdo la prima, \textbf{\$5.000}
        (100\% de lo invertido).
  \item Punto de equilibrio: \$1.150 (strike + prima).
\end{itemize}
\textbf{Ejemplo de PUT como seguro.} Tengo 100 acciones a \$1.000 y temo una caída. Compro un put
con strike \$950 pagando \$40 por acción (\$4.000). Si la acción cae a \$800, puedo vender a
\$950: el put me paga $(950 - 800 - 40) \times 100 =$ \textbf{\$11.000}, que compensa buena
parte de la caída. Si la acción sube, pierdo solo los \$4.000, como quien paga un seguro.
\end{recuadro}

\paragraph*{Para qué se usan las opciones.} Para \textbf{especular} con poco dinero
(apalancamiento), para \textbf{cubrirse} de caídas (puts como seguro) o para \textbf{generar
ingresos} lanzando opciones sobre acciones que ya se tienen (\textit{lanzamiento cubierto}).
\textbf{Advertencia:} la mayoría de las opciones compradas vencen sin valor, y el lanzador
descubierto puede tener pérdidas muy grandes. Son instrumentos para inversores con experiencia.

\subsubsection*{4.9. La venta en corto (\textit{short})}
\textbf{Ir ``corto''} es apostar a que un precio \textbf{va a bajar}. Como no se puede vender lo
que no se tiene, el inversor \textbf{toma prestadas} las acciones (préstamo de valores, por el
que paga una tasa), las vende hoy y más adelante las recompra para devolverlas.

\begin{recuadro}
\textbf{Ejemplo:} tomo prestadas 100 acciones que valen \$1.000 y las vendo: recibo \$100.000.
\begin{itemize}[leftmargin=*,nosep]
  \item Si bajan a \$800, las recompro por \$80.000 y las devuelvo: gano \textbf{\$20.000}
        (menos el costo del préstamo).
  \item Si suben a \$1.400, tengo que recomprarlas por \$140.000: pierdo \textbf{\$40.000}.
\end{itemize}
Comprando acciones, lo máximo que se puede perder es lo invertido (si la acción vale cero). En
un corto la pérdida \textbf{no tiene límite}, porque el precio puede subir indefinidamente.
\end{recuadro}

Lo contrario del corto es la \textbf{posición larga} (\textit{long}): comprar esperando que suba.
En Argentina la venta en corto está habilitada en BYMA, con préstamo de valores y garantías.

\subsubsection*{4.10. Apalancamiento y \textit{margin call}}
\textbf{Apalancarse} es invertir con dinero prestado (por ejemplo, con una caución tomadora o
con derivados) para multiplicar el resultado. Multiplica las ganancias, pero también las
pérdidas. Si las garantías pierden valor, el broker exige reponerlas (\textbf{llamado de
margen} o \textit{margin call}); si no se reponen, \textbf{vende las posiciones} aunque sea con
pérdida.

\subsubsection*{4.11. Arbitraje}
Aprovechar una diferencia de precio del mismo activo en dos mercados o monedas (por ejemplo, un
CEDEAR que en pesos quedó desalineado respecto de su precio en Nueva York). Los arbitrajistas
hacen que los precios se mantengan alineados.

% ============================================================
\section*{5. MERCADO ALCISTA Y BAJISTA: TOROS Y OSOS}

\begin{center}
\small
\begin{tabular}{|p{0.44\linewidth}|p{0.44\linewidth}|}
\hline
\textbf{Mercado alcista (\textit{bull market})} & \textbf{Mercado bajista (\textit{bear market})}
\\ \hline
El \textbf{toro} ataca de abajo hacia arriba con sus cuernos: los precios \textbf{suben}. &
El \textbf{oso} ataca de arriba hacia abajo con sus garras: los precios \textbf{bajan}. \\
\hline
Por convención, una suba de 20\% o más desde un mínimo reciente. & Por convención, una caída de
20\% o más desde un máximo reciente. \\ \hline
Optimismo, economía en crecimiento, empresas que ganan más, tasas bajas. & Pesimismo, recesión,
crisis, suba de tasas, guerras. \\ \hline
Un inversor \textbf{\textit{bullish}} (alcista) espera subas: compra o va largo. & Un inversor
\textbf{\textit{bearish}} (bajista) espera caídas: vende, se cubre con puts o va corto. \\
\hline
\end{tabular}
\end{center}

\paragraph*{Otros términos del ciclo.}
\begin{itemize}
  \item \textbf{Rally:} suba fuerte en poco tiempo.
  \item \textbf{Corrección:} caída de entre 10\% y 20\% dentro de una tendencia alcista; es
        normal y frecuente.
  \item \textbf{Crash:} desplome brusco en pocos días (1929, 1987, 2008, marzo de 2020).
  \item \textbf{Burbuja:} los precios suben muy por encima del valor real porque todos compran
        esperando que siga subiendo (las ``puntocom'' en 2000, la vivienda en EE.\,UU. en 2007).
        Cuando ``explota'', la caída es muy fuerte.
  \item \textbf{Volatilidad:} cuánto y qué tan rápido se mueven los precios. El índice
        \textbf{VIX} mide la volatilidad esperada en EE.\,UU. y se lo llama ``índice del miedo''.
  \item \textbf{Toma de ganancias:} ventas después de una suba para asegurar lo ganado.
\end{itemize}

\paragraph*{La psicología del inversor.}
Los errores más comunes no son técnicos sino emocionales: comprar por miedo a quedarse afuera
(\textbf{FOMO}) cuando algo ya subió mucho, y vender en \textbf{pánico} cuando cae. Por eso se
dice que hay que ``comprar cuando hay miedo y ser prudente cuando hay euforia''.

% ============================================================
\section*{6. CÓMO SE ANALIZA UNA INVERSIÓN}

\subsubsection*{6.1. Análisis fundamental}
Estudia el \textbf{valor real} de una empresa: sus ventas, ganancias, deudas, la calidad de sus
directivos, su sector y su futuro. Algunos indicadores:
\begin{itemize}
  \item \textbf{Capitalización bursátil:} precio de la acción $\times$ cantidad de acciones. Es
        lo que ``vale'' la empresa en la bolsa.
  \item \textbf{Ganancia por acción (EPS):} ganancia neta $\div$ cantidad de acciones.
  \item \textbf{Relación precio/ganancia (P/E o PER):} precio $\div$ ganancia por acción. Indica
        cuántos años de ganancias se pagan; un P/E muy alto puede indicar una acción cara (o
        mucha expectativa de crecimiento).
  \item \textbf{Rendimiento por dividendo:} dividendo anual $\div$ precio.
  \item \textbf{Endeudamiento:} deuda en relación al patrimonio.
\end{itemize}

\subsubsection*{6.2. Análisis técnico}
Estudia los \textbf{gráficos de precios y volúmenes} buscando patrones:
\begin{itemize}
  \item \textbf{Velas japonesas:} cada vela muestra el precio de apertura, cierre, máximo y mínimo
        de un período.
  \item \textbf{Tendencias:} alcista (máximos y mínimos cada vez más altos), bajista o lateral.
  \item \textbf{Soportes y resistencias:} niveles donde el precio suele dejar de bajar o de subir.
  \item \textbf{Medias móviles:} promedios del precio de los últimos 50 o 200 días, para
        suavizar el ruido diario.
\end{itemize}
Ningún análisis predice el futuro con seguridad. Los profesionales suelen combinar ambos y,
sobre todo, administrar el riesgo.

% ============================================================
\section*{7. LAS BOLSAS DEL MERCOSUR Y DE LA REGIÓN}

\begin{center}
\small
\begin{tabular}{|p{0.20\linewidth}|p{0.27\linewidth}|p{0.43\linewidth}|}
\hline
\textbf{País} & \textbf{Bolsa} & \textbf{Características} \\ \hline
Brasil & \textbf{B3} (Brasil, Bolsa, Balcão), San Pablo & La bolsa más grande de América Latina.
Su índice es el \textbf{Ibovespa}, con empresas como Petrobras, Vale, Itaú y Embraer. \\ \hline
Uruguay & Bolsa de Valores de Montevideo & Mercado chico, centrado en bonos del Estado
uruguayo. \\ \hline
Paraguay & Bolsa de Valores de Asunción & Mercado en crecimiento, principalmente de renta fija.
\\ \hline
Chile (no Mercosur) & Bolsa de Santiago & Índice IPSA. Uno de los mercados más desarrollados de
la región. \\ \hline
México (no Mercosur) & Bolsa Mexicana de Valores & Índice S\&P/BMV IPC. \\ \hline
\end{tabular}
\end{center}

% ============================================================
\section*{8. LAS BOLSAS DE ESTADOS UNIDOS Y DEL MUNDO}

\subsubsection*{8.1. Estados Unidos: el mercado más grande del mundo}
\begin{itemize}
  \item \textbf{NYSE} (New York Stock Exchange): la bolsa de Wall Street, la más grande del
        mundo por valor de las empresas listadas.
  \item \textbf{Nasdaq:} bolsa electrónica donde cotizan la mayoría de las grandes empresas
        tecnológicas (Apple, Microsoft, Nvidia, Alphabet, Amazon, Meta).
  \item \textbf{ADR:} certificados con los que empresas extranjeras cotizan en EE.\,UU. Varias
        argentinas lo hacen (YPF, Grupo Galicia, Pampa Energía, Telecom); su precio en Nueva
        York es una referencia de cómo ve el mundo a la Argentina.
\end{itemize}

\begin{center}
\small
\begin{tabular}{|p{0.22\linewidth}|p{0.68\linewidth}|}
\hline
\textbf{Índice} & \textbf{Qué mide} \\ \hline
\textbf{Dow Jones Industrial Average} & El más antiguo (1896). Solo 30 grandes empresas
``tradicionales'' (Boeing, Coca-Cola, JPMorgan, McDonald's...). Se calcula por precio. \\ \hline
\textbf{S\&P 500} & Las 500 empresas más grandes de EE.\,UU. Es el índice de referencia para medir
``al mercado'' estadounidense. \\ \hline
\textbf{Nasdaq Composite / Nasdaq 100} & Todas las empresas del Nasdaq (o las 100 mayores no
financieras). Muy concentrado en tecnología. \\ \hline
\textbf{Russell 2000} & Empresas pequeñas. \\ \hline
\end{tabular}
\end{center}

\subsubsection*{8.2. Otros mercados del mundo}
\begin{center}
\small
\begin{tabular}{|p{0.22\linewidth}|p{0.30\linewidth}|p{0.38\linewidth}|}
\hline
\textbf{Región} & \textbf{Bolsa} & \textbf{Índice principal} \\ \hline
Europa & Euronext, Fráncfort, Londres & Euro Stoxx 50, DAX (Alemania), CAC 40 (Francia),
FTSE 100 (Reino Unido), FTSE MIB (Italia) \\ \hline
Japón & Bolsa de Tokio & Nikkei 225 \\ \hline
China & Shanghái, Shenzhen, Hong Kong & Shanghai Composite, Hang Seng \\ \hline
Global & --- & MSCI World, MSCI Emerging Markets (índices que reúnen muchos países) \\ \hline
\end{tabular}
\end{center}

\paragraph*{El mundo conectado.} Los mercados se influyen entre sí: una suba de tasas en
Estados Unidos, un conflicto armado o una crisis en China mueven los precios de acciones,
bonos y \textbf{materias primas} (petróleo, oro, soja, cobre) en todo el mundo, incluida la
Argentina. Para un país exportador de granos y energía, el precio internacional de la soja y
del petróleo es tan importante como el Merval.

\subsubsection*{8.3. Bonos del mundo}
\begin{itemize}
  \item \textbf{Bonos del Tesoro de EE.\,UU. (Treasuries):} considerados la inversión más segura
        del mundo; su rendimiento a 10 años es la base para medir el riesgo de los demás países.
  \item \textbf{Bunds} (Alemania), \textbf{Gilts} (Reino Unido), \textbf{JGB} (Japón): bonos de
        países desarrollados, también de bajo riesgo.
  \item \textbf{Bonos de países emergentes} (Argentina, Brasil, Turquía...): pagan más porque
        tienen más riesgo.
  \item Las \textbf{calificadoras de riesgo} (S\&P, Moody's, Fitch) ponen notas a países y
        empresas: AAA es la máxima; desde BB+ hacia abajo se considera ``especulativo''.
\end{itemize}

% ============================================================
\section*{9. INSTRUMENTOS PARA EL PEQUEÑO INVERSOR}

\subsubsection*{9.1. Fondos comunes de inversión (FCI)}
Un \textbf{FCI} junta el dinero de muchos ahorristas y lo invierte un profesional (la sociedad
gerente) según un reglamento aprobado por la CNV. Cada ahorrista recibe
\textbf{cuotapartes}.
\begin{itemize}
  \item \textbf{Money market:} invierten en plazos fijos y cuentas remuneradas; rescate
        inmediato y bajo riesgo. Son los que usan muchas billeteras virtuales para dar
        rendimiento diario.
  \item \textbf{Renta fija:} bonos y letras, en pesos o en dólares.
  \item \textbf{Renta variable:} acciones.
  \item \textbf{Mixtos}, \textbf{pymes} e \textbf{infraestructura}.
\end{itemize}
\textbf{Ventajas:} diversificación con poco dinero y gestión profesional.
\textbf{Contras:} cobran honorarios y, salvo los money market, su valor puede bajar.

\subsubsection*{9.2. ETF (Exchange Traded Funds)}
Un \textbf{ETF} es un fondo que \textbf{cotiza en la bolsa como si fuera una acción} y que, en
general, \textbf{replica un índice}. Por ejemplo, el ETF \textbf{SPY} o el \textbf{IVV}
replican el S\&P 500: comprando una sola unidad se tiene una pequeña parte de las 500
empresas. Hay ETF de países, sectores (tecnología, energía), materias primas (oro, cobre) y
bonos. Tienen comisiones bajas y son una forma sencilla de diversificar. La inversión
\textbf{pasiva} (comprar el índice y mantenerlo muchos años) suele superar, a largo plazo, a la
mayoría de los inversores que intentan ``ganarle al mercado''.

\subsubsection*{9.3. CEDEARs}
Los \textbf{CEDEARs} (Certificados de Depósito Argentinos) permiten comprar en BYMA, \textbf{en
pesos o dólares}, acciones y ETF que cotizan en el exterior (Apple, Coca-Cola, Mercado Libre,
Boeing, el SPY). Un banco o la Caja de Valores compra el título original en el exterior y
emite certificados en la Argentina con un \textbf{ratio} (por ejemplo, 10 CEDEARs = 1 acción).
\begin{itemize}
  \item Como el activo está en dólares, su precio en pesos sigue al dólar
        (\textbf{cobertura cambiaria}) y además a la acción.
  \item Hoy hay \textbf{más de 400 CEDEARs} disponibles en BYMA, incluidos muchos de ETF.
  \item Riesgos: los de la empresa o índice, más el tipo de cambio.
\end{itemize}

\subsubsection*{9.4. Criptoactivos}
Bitcoin, Ethereum y las \textbf{stablecoins} (como USDT o USDC, que siguen al dólar) se compran
en exchanges. En Argentina, las plataformas deben estar inscriptas en la CNV como
\textbf{Proveedores de Servicios de Activos Virtuales (PSAV)}. Son \textbf{muy volátiles}, no
tienen respaldo de un Estado ni de una empresa, y son el terreno preferido de muchas estafas.

\subsubsection*{9.5. Comparación rápida}
\begin{center}
\footnotesize
\setlength{\tabcolsep}{3pt}
\begin{tabular}{|>{\raggedright\arraybackslash}p{0.17\linewidth}|>{\raggedright\arraybackslash}p{0.17\linewidth}|>{\raggedright\arraybackslash}p{0.15\linewidth}|>{\raggedright\arraybackslash}p{0.17\linewidth}|>{\raggedright\arraybackslash}p{0.22\linewidth}|}
\hline
\textbf{Instrumento} & \textbf{Qué compro} & \textbf{Riesgo} & \textbf{Liquidez} &
\textbf{Para qué sirve} \\ \hline
Plazo fijo & Un depósito bancario & Bajo & Al vencimiento & Ahorro de corto plazo en pesos \\
\hline
FCI money market & Cuotapartes de un fondo & Bajo & Inmediata & Tener el dinero disponible y
rindiendo \\ \hline
Caución colocadora & Un préstamo de días & Bajo & A los pocos días & Colocar pesos o dólares
de corto plazo \\ \hline
Bonos y ON & Deuda de un Estado o empresa & Bajo a alto, según emisor & Alta (bolsa) & Renta
periódica \\ \hline
Acciones & Parte de una empresa & Alto & Alta (bolsa) & Largo plazo \\ \hline
ETF / CEDEARs & Una canasta de acciones o una acción del exterior & Medio a alto & Alta (bolsa)
& Diversificar, largo plazo \\ \hline
Futuros y opciones & Un contrato derivado & Muy alto & Alta & Cobertura o especulación \\
\hline
Criptoactivos & Un activo digital & Muy alto & Alta & Especulación \\ \hline
\end{tabular}
\end{center}

% ============================================================
\section*{10. LOS BANCOS CENTRALES: BCRA Y RESERVA FEDERAL}

Un \textbf{banco central} es la institución que emite la moneda de un país y regula el sistema
financiero. Sus decisiones mueven todos los mercados.

\subsubsection*{10.1. Banco Central de la República Argentina (BCRA)}
\begin{itemize}
  \item \textbf{Emite} el peso y administra la cantidad de dinero en la economía.
  \item Busca la \textbf{estabilidad monetaria y financiera} (cuidar el valor del peso).
  \item Administra las \textbf{reservas internacionales} (dólares, oro y otras divisas que
        respaldan la moneda y permiten pagar importaciones y deuda).
  \item Fija la \textbf{política cambiaria} (cómo se determina el precio del dólar oficial).
  \item \textbf{Regula y supervisa} a bancos y entidades financieras y protege a los usuarios.
\end{itemize}

\subsubsection*{10.2. Reserva Federal de Estados Unidos (Fed)}
La \textbf{Fed} es el banco central de Estados Unidos y emite el dólar, la moneda de reserva
del mundo. Tiene un \textbf{doble mandato}: \textbf{pleno empleo} y \textbf{estabilidad de
precios} (inflación cercana al 2\% anual). Su principal herramienta es la \textbf{tasa de
fondos federales}, que decide el Comité Federal de Mercado Abierto (FOMC) en ocho reuniones al
año.

\begin{recuadro}
\textbf{¿Por qué a la Argentina le importa la tasa de la Fed?}
\begin{itemize}[leftmargin=*,nosep]
  \item Si la Fed \textbf{sube} la tasa, los inversores pueden ganar más en bonos de EE.\,UU.
        sin riesgo, entonces se van de los países emergentes: suben el riesgo país y el dólar,
        bajan las acciones y los bonos argentinos, y se encarece la deuda.
  \item Si la Fed \textbf{baja} la tasa, ocurre lo contrario: llegan más inversiones a los
        emergentes.
\end{itemize}
\textbf{Dato:} en su reunión del 28 de octubre de 2026 la Fed mantuvo la tasa en el rango de
3,75\% a 4\%, preocupada por una inflación que sigue alta.
\end{recuadro}

% ============================================================
\section*{11. EL TIEMPO COMO ALIADO: INTERÉS COMPUESTO}

El \textbf{interés compuesto} es ganar intereses sobre los intereses ya ganados. Con el tiempo
el efecto se vuelve enorme.

\begin{recuadro}
\textbf{Ejemplo:} ahorrar \textbf{USD 100 por mes} en una inversión que rinde en promedio 7\%
anual.
\begin{itemize}[leftmargin=*,nosep]
  \item En \textbf{10 años} se aportaron USD 12.000 y se acumulan unos \textbf{USD 17.300}.
  \item En \textbf{30 años} se aportaron USD 36.000 y se acumulan unos \textbf{USD 122.000}.
\end{itemize}
Empezar joven es la mayor ventaja que tiene un inversor. (El 7\% es un supuesto de ejemplo; los
rendimientos reales varían año a año y pueden ser negativos.)
\end{recuadro}

% ============================================================
\section*{12. CUIDADOS DEL INVERSOR}

\begin{itemize}
  \item \textbf{Primero, un fondo de emergencia:} antes de invertir en algo riesgoso, tener
        ahorrados algunos meses de gastos en algo líquido y seguro.
  \item \textbf{Invertir solo dinero que no se necesita} en el corto plazo.
  \item \textbf{Entender lo que se compra:} si no se puede explicar en dos frases, mejor no.
  \item \textbf{Diversificar} entre instrumentos, monedas y plazos.
  \item \textbf{Operar solo con agentes registrados en la CNV}: se puede verificar en su sitio.
  \item \textbf{Desconfiar} de quien promete ganancias altas y seguras, ``esquemas'' de
        referidos, ``trading'' guiado por grupos de WhatsApp o Telegram, o supuestos expertos
        en redes sociales: son las estafas financieras más comunes (esquemas Ponzi).
  \item \textbf{Nunca} compartir claves ni códigos de la cuenta del broker.
  \item \textbf{Informarse en fuentes confiables:} CNV, BCRA, BYMA Educa, prensa económica.
\end{itemize}

% ============================================================
\section*{13. GLOSARIO BURSÁTIL}
\begin{center}
\footnotesize
\begin{tabular}{|p{0.22\linewidth}|p{0.68\linewidth}|}
\hline
\textbf{Término} & \textbf{Significado} \\ \hline
Acción & Parte del capital de una empresa. \\ \hline
\textit{Blue chip} & Acción de una empresa grande, sólida y muy negociada. \\ \hline
Boleto & Comprobante de cada operación realizada. \\ \hline
Cartera (portafolio) & Conjunto de inversiones de una persona. \\ \hline
Cotización & Precio de un título en el mercado. \\ \hline
Cupón & Pago de intereses de un bono. \\ \hline
Dividendo & Ganancia que una empresa reparte entre sus accionistas. \\ \hline
\textit{Hedge} (cobertura) & Operación para protegerse de un riesgo. \\ \hline
IPO & Oferta pública inicial: cuando una empresa sale a cotizar por primera vez. \\ \hline
Liquidez & Facilidad para vender un activo rápido sin perder precio. \\ \hline
Long / Short & Posición compradora (gana si sube) / vendedora (gana si baja). \\ \hline
Prima & Precio que se paga por una opción. \\ \hline
Rueda & Jornada de negociación. \\ \hline
Spread & Diferencia entre la punta compradora y la vendedora. \\ \hline
Strike & Precio de ejercicio de una opción. \\ \hline
Tenencia & Lo que tiene el inversor en su cuenta. \\ \hline
Ticker & Código abreviado de un título. \\ \hline
Volatilidad & Intensidad con que varían los precios. \\ \hline
\end{tabular}
\end{center}

% ============================================================
\section*{14. ACTIVIDADES DEL MÓDULO}

\begin{enumerate}
  \item \textbf{Termómetro del mercado:} durante dos semanas, registrar en una planilla el
        Merval, el riesgo país, el dólar oficial y MEP, el S\&P 500, el Ibovespa y el precio de
        un bono (por ejemplo, GD30). Graficar y relacionar los movimientos con las noticias.
        ¿El mercado está \textit{bullish} o \textit{bearish}?
  \item \textbf{Riesgo país:} con el rendimiento del bono del Tesoro de EE.\,UU. a 10 años y el
        riesgo país del día, calcular cuánto rinden aproximadamente los bonos argentinos en
        dólares.
  \item \textbf{Comparar brokers:} elegir tres brokers o bancos y comparar requisitos de apertura,
        comisiones, instrumentos disponibles y facilidad de uso. Verificar su registro en la CNV.
  \item \textbf{Dólar MEP:} con las cotizaciones del día de AL30 y AL30D (o GD30 y GD30D),
        calcular el dólar MEP y compararlo con el oficial.
  \item \textbf{Cartera simulada:} con \$1.000.000 ficticios, armar una cartera diversificada
        (FCI money market, caución, un bono, acciones argentinas y CEDEARs) y seguirla durante
        un mes. Comparar con haber comprado dólares o hecho un plazo fijo.
  \item \textbf{Opciones en papel:} con el precio actual de GGAL y las primas de sus opciones,
        calcular la ganancia o pérdida de comprar un call y un put si al vencimiento la acción
        sube 15\% o baja 15\%.
  \item \textbf{Cobertura con futuros:} un taller debe pagar USD 5.000 en dos meses; con el precio
        del dólar futuro del día, explicar cómo se cubriría y qué pasa si el dólar sube o baja.
  \item \textbf{Investigar una empresa:} elegir una empresa que cotice en BYMA o un CEDEAR del
        sector aeronáutico o tecnológico (por ejemplo, Boeing o Mercado Libre; o Embraer en la
        B3). Explicar qué hace, en qué bolsa cotiza, su capitalización, su P/E y cómo evolucionó
        su precio.
  \item \textbf{La Fed y nosotros:} buscar la última decisión de la Fed y del BCRA y explicar
        cómo pueden afectar el dólar, el riesgo país y el precio de los bonos argentinos.
  \item \textbf{Detectar la estafa:} analizar publicidades de ``inversiones'' en redes sociales
        e identificar señales de alerta.
\end{enumerate}
